\documentclass[12pt]{article}

%\geometry{showframe}% for debugging purposes -- displays the margins



\newcommand{\bra}[1]{\left(#1\right)}
\usepackage{clrscode3e}
\usepackage{xcolor}
\usepackage{tikz}
\usepackage{amsmath,amsthm}
\usepackage{amssymb}
\usetikzlibrary{shapes}
\usetikzlibrary{positioning}
% Set up the images/graphics package
\usepackage{graphicx}
\setkeys{Gin}{width=\linewidth,totalheight=\textheight,keepaspectratio}
\graphicspath{{graphics/}}

\title{DRP write up}
\author{Shan Gao}
\date{Winter 2024}  % if the \date{} command is left out, the current date will be used

% The following package makes prettier tables.  We're all about the bling!
\usepackage{booktabs}

% The units package provides nice, non-stacked fractions and better spacing
% for units.
\usepackage{units}

% The fancyvrb package lets us customize the formatting of verbatim
% environments.  We use a slightly smaller font.
\usepackage{fancyvrb}
\fvset{fontsize=\normalsize}

% Small sections of multiple columns
\usepackage{multicol}

%--------Theorem Environments--------
%theoremstyle{plain} --- default
\newtheorem{thm}{Theorem}
\newtheorem{cor}[thm]{Corollary}
\newtheorem{prop}[thm]{Proposition}
\newtheorem{lem}[thm]{Lemma}
\newtheorem{conj}[thm]{Conjecture}
\newtheorem{quest}[thm]{Question}

\theoremstyle{definition}
\newtheorem{defn}[thm]{Definition}
\newtheorem{defns}[thm]{Definitions}
\newtheorem{con}[thm]{Construction}
\newtheorem{exmp}[thm]{Example}
\newtheorem{exmps}[thm]{Examples}
\newtheorem{notn}[thm]{Notation}
\newtheorem{notns}[thm]{Notations}
\newtheorem{addm}[thm]{Addendum}
\newtheorem{exer}[thm]{Exercise}

\theoremstyle{remark}
\newtheorem{rem}[thm]{Remark}
\newtheorem{rems}[thm]{Remarks}
\newtheorem{warn}[thm]{Warning}
\newtheorem{sch}[thm]{Scholium}


% These commands are used to pretty-print LaTeX commands
\newcommand{\doccmd}[1]{\texttt{\textbackslash#1}}% command name -- adds backslash automatically
\newcommand{\docopt}[1]{\ensuremath{\langle}\textrm{\textit{#1}}\ensuremath{\rangle}}% optional command argument
\newcommand{\docarg}[1]{\textrm{\textit{#1}}}% (required) command argument
\newenvironment{docspec}{\begin{quote}\noindent}{\end{quote}}% command specification environment
\newcommand{\docenv}[1]{\textsf{#1}}% environment name
\newcommand{\docpkg}[1]{\texttt{#1}}% package name
\newcommand{\doccls}[1]{\texttt{#1}}% document class name
\newcommand{\docclsopt}[1]{\texttt{#1}}% document class option name

\begin{document}

\maketitle
\begin{abstract}
    Presents some results found in \cite{devroye2015measure}
\end{abstract}

\section{Introduction}
\begin{thm}[Absolutely continuous w.r.t Lebesgue measure]
    We say that a measure $\mu$ on $\mathcal{B}(\mathbb{R}^d)$ is absolutely continuous with respect to the Lebesgue measure $\lambda$ if for all $A \in \mathcal{B}(\mathbb{R}^d)$ and $m(A) = 0$ then $\mu(A) =0$. 
\end{thm}
\begin{cor}
    Let $\mu$ be absolutely continuous with respect to the Lebesgue measure, then for any $A \in \mathcal{B}(\mathbb{R}^d)$, denoting the density function of $\mu$ by $f$ we have 
    $$
        \mu(A) = \int_A f dx
    $$
\end{cor}
\begin{defn}[Random Voronoi Tesselation]
        Let $X_1,\ldots,X_n$ be i.i.d random vectors taking values in $\mathbb{R}^d$. That means all of them share a common distribution , call it $\mu$. We assume that $\mu$ is absolutely continuous with respect to the Lebesgue measure $\lambda$, and denote the density function of $\mu$ by $f$. These $X_i$'s define a partition of $\mathbb{R}^d$ into $n$ sets $S_1,\ldots,S_n$ such that $S_i$ contain all points in $\mathbb{R}^d$ whose distances to $X_i$ is not greater than their distances to all other $X_j$'s ($ j \neq i $). Letting $\|.\|$ denote the standard Euclidean norm, we have that 
        $$ S_i = \{x \in \mathbb{R}^d \colon \|x - X_i\| = \min_{j = [n]} \|x - X_j\| \}   \cap \{x \in \mathbb{R}^d \colon \|x - X_i\| < \min_{j = 1,\ldots,i-1} \|x - X_j\|\}$$
\end{defn}

\begin{prop}[Asymptotic second moment]
	The asymptotic second moment is expressed in terms of a random variable $W$ defined as follows. Let $Y$ be a random vector uniformly distributed in $B(0,1)$. Define, $\vec{1} = (1,0,0,\ldots,0) \in \mathbb{R}^d$ and let $\overline{B} = B(\vec{1},1)\cup B(Y, \|Y\|)$.  \textcolor{red}{Is this assuming standard Euclidean norm?}	
	
\end{prop}
We introduce the random variable
$$
		W = \frac{m(\overline{B})}{m(B(0,1))}
$$
and let
$$
	\alpha(d) \equiv \mathbb{E}\left[ \frac{2}{W^2} \right]
$$

\begin{thm}
	Assume $\mu$ has density $f$. Then 
	$$
		n \mathbb{E} [ \mu(S_1) \colon X_1 = x] \to 1
	$$
	for $\mu$ almost all $x$.
\end{thm}





\begin{proof}
	Observe that  

	\begin{align*}
		\mathbb{E}(\mu(S_1) \colon X_1 = x) &= \mathbb{P}\{X_{n+1} \in S_1 \colon X_1 = x\}\\
	\end{align*}
	Where we used a neat trick to interchange $\mu(S_1)$ with $ X_{n+1} \in S_1$ since $X_{n+1}$ is $\mu$ distributed the probability of it landing in $S_1$ is the same as the $\mu$ measure of $S_1$. 
	Furthermore we have that 
	$$
		\mathbb{P}(X_{n+1} \in S_1 \colon X_1 = x) = \mathbb{P}\left( \bigcap_{i=2}^n \left\{ X_i \notin B\left(X_{n+1}, \|X_{n+1}- x\|\right)\right\}\right)
	$$
	This is because since $X_{n+1} \in S_1$ so it must be at least $\|X_{n+1} - x\|$ distance away from any other $X_i$ $(i \geq 2)$. Defining $Z(x) = \mu(B(X,\|X-x\|))$ and since all the $X_i$'s are indenpendently distributed, we get that 

	\begin{align*}
		\mathbb{P}\left( \bigcap_{i=2}^n \left\{ X_i \notin B\left(X_{n+1}, \|X_{n+1}- x\|\right)\right\}\right) &= \mathbb{E}[(1-Z(x))^{n-1}]
	\end{align*}
	
\end{proof}


$Z(x) = \mu(B(X,\|X-x\|))$






\bibliographystyle{plain}
\bibliography{writeup}

\end{document}




